\section{Introdução}\label{sec:introducao}

A coleta de resíduos sólidos urbanos (RSU) é um serviço contínuo, caro e muito
visível para a população. Os indicadores nacionais de 2024 registram, para o
Brasil, 1,14 kg de RSU coletado por habitante por dia e, considerando
apenas os resíduos domiciliares, 0,84 kg/hab$\cdot$dia; no Nordeste,
esses valores são 1,21 e 0,88 kg/hab$\cdot$dia, respectivamente
\cite[p.~25--26]{brasil2025sinisa}. Em uma capital como Aracaju, com
602.757 habitantes no Censo de 2022 \cite{ibge2022censo}, isso equivale a
centenas de toneladas de resíduos domiciliares por dia, que precisam ser
recolhidas por frotas com capacidade, jornada e custo limitados. A Política
Nacional de Resíduos Sólidos \cite{brasil2010lei} e o marco legal do saneamento
\cite{brasil2020lei} reforçam a exigência de planejamento e de indicadores para
esse serviço.

Na prática mais comum, a coleta segue roteiros e frequências fixos: cada
contêiner é visitado em dias predeterminados, esteja cheio ou quase vazio. Essa
regra é simples de operar, mas não se ajusta às variações de geração entre
locais e entre dias. Contêineres que enchem rápido transbordam antes da visita;
outros são esvaziados com pouca carga. A Internet das Coisas (IoT) oferece uma
alternativa: sensores de nível instalados nos contêineres transmitem a ocupação
para uma plataforma, que decide quais pontos visitar e em que ordem
\cite{anagnostopoulos2017challenges,esmaeilian2018future}. A literatura relata
reduções de distância percorrida com esse tipo de estratégia, mas também mostra
que os ganhos estimados em simulação tendem a ser maiores que os observados em
implantações reais. Uma meta-análise recente encontrou redução média de
39,79\% nos estudos de simulação e de 12,37\% nos estudos de campo
\cite{maciel2025impact}. Parte dessa diferença pode decorrer de hipóteses
favoráveis demais nos modelos, como sensores sem erro, comunicação perfeita e
comparações com linhas de base pouco exigentes.

Este trabalho é a continuação do Trabalho de Conclusão de Curso apresentado em
2024\footnote{A versão de 2024 está preservada sem alterações no repositório do
projeto, no diretório \texttt{original/}.}, no qual foi montado, a partir da
arquitetura em três camadas de \cite{oliveira2018arquitetura}, um protótipo com ESP32,
sensor ultrassônico HC-SR04, sensor de temperatura e umidade DHT11 e
comunicação MQTT. O protótipo demonstrou a troca de mensagens entre dispositivo,
broker e aplicativo, mas não avaliou se o monitoramento melhoraria a coleta, e
o firmware continha erros que invertiam o significado da medida de enchimento
(Seção~\ref{sec:firmware}). A presente versão mantém o tema e o hardware
original, corrige o firmware e acrescenta o que faltava para responder à
pergunta operacional: uma plataforma de simulação em que contêineres virtuais
publicam telemetria por MQTT, uma camada de ingestão que estima o estado de cada
contêiner e um módulo de planejamento que compara estratégias de coleta sob as
mesmas condições.

\subsection{Problema e questões de pesquisa}

O problema investigado é \emph{como uma arquitetura IoT com sensores de nível,
comunicação MQTT e planejamento orientado por dados pode apoiar a eficiência
da coleta de RSU, considerando as imperfeições de medição e de comunicação}.
Ele foi desdobrado em cinco questões:

\begin{description}
  \item[QP1] Como o monitoramento do nível de enchimento pode apoiar a decisão
  de quando coletar?
  \item[QP2] Em cenários equivalentes, que diferenças operacionais existem entre
  a coleta periódica fixa, a coleta por limiar de enchimento e a coleta
  dinâmica com otimização de rotas, com ou sem previsão?
  \item[QP3] Como perda e atraso de mensagens, indisponibilidade de sensores e
  ruído de medição afetam essas estratégias?
  \item[QP4] Qual é a influência da quantidade e da distribuição dos
  contêineres, da demanda e da frota sobre os resultados?
  \item[QP5] Em que condições a estratégia orientada por dados traz ganhos
  relevantes e em quais a complexidade adicional não se justifica?
\end{description}

\subsection{Objetivos}

O objetivo geral é desenvolver e avaliar experimentalmente, por simulação, uma
plataforma IoT baseada em MQTT para monitoramento de contêineres e planejamento
da coleta de RSU. Os objetivos específicos são:
(i)~auditar e corrigir o firmware do protótipo de 2024;
(ii)~especificar um contrato de mensagens MQTT com identificação de
dispositivo, sequência, validade e versão de esquema;
(iii)~implementar um simulador de contêineres, sensores ultrassônicos e canal de
comunicação com injeção de falhas;
(iv)~implementar a ingestão, a estimação de estado e quatro políticas de coleta
com roteirização por otimização;
(v)~executar experimentos pareados e reprodutíveis em um cenário construído
sobre a malha viária real de Aracaju; e
(vi)~analisar os resultados com estatística apropriada e discutir limites de
validade.

\subsection{Delimitação e contribuições}

Nenhum componente físico novo foi montado e nenhum dado de campo foi coletado.
Toda a avaliação quantitativa é feita por simulação; a malha viária é real
(OpenStreetMap), mas as posições dos contêineres e do depósito são
hipotéticas, e os resultados não representam o serviço de coleta de Aracaju.
Os parâmetros de geração de resíduos combinam dados oficiais e hipóteses
declaradas (Seção~\ref{sec:parametros}).

As contribuições efetivamente demonstradas neste trabalho são:
(a)~a auditoria e a reescrita do firmware, com compilação verificada;
(b)~um contrato MQTT~3.1.1 e uma camada de ingestão que trata duplicatas,
mensagens fora de ordem, leituras inválidas e dados desatualizados;
(c)~uma plataforma de simulação aberta, com testes automatizados e verificação
de balanço de volume, que usa o mesmo código de ingestão tanto com o broker
real quanto com um canal emulado;
(d)~um protocolo experimental pareado (números aleatórios comuns) com
\NExecTotal{} execuções de simulação e testes de integração com o broker; e
(e)~evidência quantitativa sobre a troca entre distância percorrida e nível de
serviço, e sobre a sensibilidade dessas estratégias a falhas de telemetria.

\subsection{Organização do texto}

A Seção~\ref{sec:fundamentacao} apresenta os conceitos usados no trabalho. A
Seção~\ref{sec:relacionados} discute trabalhos relacionados e posiciona a
proposta. A Seção~\ref{sec:arquitetura} descreve a arquitetura, o firmware e a
plataforma. A Seção~\ref{sec:metodologia} detalha os modelos, os parâmetros e o
protocolo experimental. A Seção~\ref{sec:resultados} apresenta e discute os
resultados, e a Seção~\ref{sec:conclusao} conclui o trabalho.
